\section{The Fragility of Civilisation}

\begin{quotationx}
\textit{Ne craignons pas de leur donner de l'amour-propre : confessions que ces ennemis de l'État, de tous les États, sont redoutables, encore que fort médiocre et peu nombreux. Car, s'il faut de long âges, un effort méthodique et persévérant, des inventions presque divines pour bâtir une ville, élever un État, constituer une civilisation, il n'y a rien de plus aisé que de défaire ces délicates compositions. Quelques tonnes de poudre vile renversent un moitié du Parthénon : une colonie de microbes décime le people d'Athènes ; trios ou quatre basses idées systématisées par des sots n'ont point mal réussi depuis un siècle à rendre vains mille ans d'Histoire de France.}

\textit{Que le mensonge libéral se propage donc sur la Terre, que l'anarchisme et le démocratisme universel étendent la pambéotie annoncée par Renan, que les barbares d'en bas, prédits par Macaulay, paraissent à propos, et l'homme pourra disparaitre en tant qu'être humain, ainsi qu'il aura disparu au titre de Français, de Grec ou de Latin.}
\end{quotationx}


\begin{quotationx}
Let us not fear giving them their proper respect: let us confess that the enemies of the State, of all States, are formidable, even if second-rate and less numerous. For, if methodical and persistent efforts for long ages, almost divine inventions, are necessary to build a city, erect a State, establish a civilisation, there is nothing easier than to undo these fragile formations. A few tons of vile powder knocks down half of the Pantheon; a colony of microbes decimates the population of Athens, three or four base ideas systematised by fools have succeeded for a century to render useless a thousand years of French history.

Should the liberal lie spread over the earth, should anarchism and universal democratism extend the general stupidity announced by Renan, should the barbarians from below predicted by Macaulay appear at the right moment, and man will disappear as a human being just as he will have disappeared as a Frenchman, a Greek, or a Latin.
\end{quotationx}


\flright{\textsc{Charles Maurras}, \emph{Quand les Français ne s'Aimaient pas}}


\flrightit{Posted on 2008-11-24 by Cologero}

